% page 460 du fac-simile (image p-459 du scan)
% bloc 154.9 x 229.0 mm ; marge gauche 8.0 mm
\pgc{-12.700mm}{0.000mm}{
\l{\cel{3}452}
\l{}
\l{\sou{pompo}. I. Procesiono solena. - II. Expozo di ulo splendida,}
\l{\cel{3}luxozega, quan onu vidigas pro ostento. DEFIS.}
\l{}
\l{\sou{pompelmuso}. (\sou{bot}\cel{1}.) Speco di \textquotedbl{}citrus\textquotedbl{} di qua la frukto,gro-}
\l{\cel{3}sega, konsumesas da la Usani e da la Britaniani lor lia}
\l{\cel{3}dejuneto. Lu kultivesas omnaloke en la regioni tropikala.}
\l{\cel{3}- L.\cel{1}\sou{citrus decumana.}\cel{1}- DFIS.}
\l{}
\l{\sou{pompono}. Tufo per qua onu ornas la kapvesti, la vestizuri}
\l{\cel{3}hominala. Tufo de lano quan la soldati +weras somite di}
\l{\cel{3}sua kepio. - DEFS.}
\l{+\sou{pondar}. (\sou{aludante ucelino}). Emisar sua ovi. - F.}
\l{\sou{ponderar}. (\sou{trans.})\cel{2}I. Mezurar la pezo di korpo, per kompa-}
\l{\cel{3}rar lu a pez-unajo. - II. (\sou{metaf.}) Judikar pri la valoro}
\l{\cel{3}di kozo, per explorar la \textquotedbl{}por\textquotedbl{} e la \textquotedbl{}kontre\textquotedbl{}. - EFIS.}
\l{}
\l{\sou{poneo.}\cel{3}(\sou{zool}.)Kavalo di raso karakterizata da staturo}
\l{\cel{3}ne-alta. - DEF.}
\l{}
\l{\sou{poniardo}. Armo atakala, kun lamo kurta, akuta, fixigita}
\l{\cel{3}an mancho, quan onu tenas per sua pugno. -\cel{2}EFIS.}
\l{}
\l{\sou{ponsa}. Koloro dazlante-reda, kelke simila a ta di papave-}
\l{\cel{3}reto. -\cel{2}DFIRS.}
\l{}
\l{\sou{ponto}. Konstrukturo quan onu uzas por pasar de ca a ta}
\l{\cel{3}rivo di rivero, di fluvio, di fosato, di granda kavajo}
\l{\cel{3}en la surfaco di sulo. - FIS.}
\l{}
\l{\sou{pontifiko.}\cel{2}I. Pastoro di religio. - II. Che la katoliki :-}
\l{\cel{3}episkopo. - DEFIRS.}
\l{}
\l{\sou{pontono}. I. Ponto flotacanta, facita ek du bateli interjun-}
\l{\cel{3}tita per trabi. - II. Konstrukturo uzata kom staciono di}
\l{\cel{3}departo o di arivo por la navi qui transportas pasajeri.}
\l{\cel{3}- DEFIRS.}
\l{}
\l{\sou{ponyeto}. Bendo a qua abutas la maniko di kamizo, di robo,di}
\l{\cel{3}korsajo, en la loko ube lu kovras la karpo. - F.}
\l{}
\l{\sou{popo.}\cel{2}Sacerdoto di la rituo Greka, che la Rusi, la Serbi}
\l{\cel{3}e la Bulgari. -}
\l{}
\l{\sou{poplo}. (\sou{bot.}) Arboro ek la familio \textquotedbl{}saliki\textquotedbl{}, kun stipo rekta,}
\l{\cel{3}e folii alternanta. L.\cel{1}\sou{populus}.-DEFI.}
\l{}
\l{\sou{poplino}. Stofo di qua la warpo-fili esas silka, e la wefto-}
\l{\cel{3}fili, lana. - DEFIRS.}
\l{}
\l{\sou{poplito}. (\sou{anat.}) Parto di la gambo, qua esas dop la genuo,}
\l{\cel{3}opoze ad ica. - ISL.}
}
